\section{The proof-producing finite-trace compiler}
\label{sec:tracecompiler}

\subsection{Trace systems}

A finite computation is internalized by fixed primitive-recursive symbols.
For clarity, consider one input parameter \(u\); tuples of parameters are
paired.

\begin{definition}[Trace system]
A trace system consists of fixed symbols
\[
\begin{array}{rclcrcl}
 \Good&:&\N\to\Bool, && \operatorname{init}&:&\N\to\N,\\
 \operatorname{bound}&:&\N\to\N, &&
 \operatorname{step}&:&\N^3\to\N,\\
 \Inv&:&\N^3\to\Bool, &&
 \Accept&:&\N^2\to\Bool.
\end{array}
\]
The run is the primitive recursion
\begin{align*}
 \Run(0,u)&=\operatorname{init}(u),\\
 \Run(Si,u)&=\operatorname{step}(i,\Run(i,u),u).
\end{align*}
\end{definition}

All predicates are represented by equations to the Boolean value one.  We use
the abbreviations
\[
\begin{aligned}
 G(u)&:\equiv \Good(u)=1,\\
 I(i,s,u)&:\equiv \Inv(i,s,u)=1,\\
 A(s,u)&:\equiv \Accept(s,u)=1.
\end{aligned}
\]

\subsection{Local proof interface}

The compiler accepts three closed object-language derivations:
\begin{align}
 \mathsf{base}:&\quad []\vdash[\PRAAx]
       G(u)\to I(0,\operatorname{init}(u),u),\label{eq:tracebase}\\
 \mathsf{step}:&\quad []\vdash[\PRAAx]
       G(u)\to I(i,s,u)\to
       I(Si,\operatorname{step}(i,s,u),u),\label{eq:tracestep}\\
 \mathsf{final}:&\quad []\vdash[\PRAAx]
       G(u)\to I(\operatorname{bound}(u),s,u)\to A(s,u).
       \label{eq:tracefinal}
\end{align}
Their displayed free variables are represented by de Bruijn indices.  The
conclusion to construct is
\begin{equation}
 []\vdash[\PRAAx]
 \forall u\,[G(u)\to
 A(\Run(\operatorname{bound}(u),u),u)].
 \label{eq:traceresult}
\end{equation}

\begin{theorem}[Finite-trace reflection]
\label{thm:tracereflection}
For every fixed trace system, the three derivations
\eqref{eq:tracebase}--\eqref{eq:tracefinal} primitive-recursively determine a
derivation of \eqref{eq:traceresult}.
\end{theorem}

\subsection{Human proof of trace reflection}

We give the proof at the same granularity as the intended Lean constructor.
No semantic induction outside iPRA is used.

\paragraph{Step 1: open the input and assume the guard.}
Work in one free variable \(u\) and context \([G(u)]\).  Because
\(\mathsf{base}\) is closed, substitute the term \(u\) for its free variable
and weaken it to this context.  Apply implication elimination to obtain
\[
    [G(u)]\vdash I(0,\operatorname{init}(u),u).
\]
The primitive-recursion zero equation gives
\[
    [G(u)]\vdash
    \Run(0,u)=\operatorname{init}(u).
\]
Use Leibniz transport in the state position of \(I(0,-,u)\) to derive
\begin{equation}
   [G(u)]\vdash I(0,\Run(0,u),u).
   \label{eq:invbase}
\end{equation}

\paragraph{Step 2: construct the induction step.}
The induction formula, with free variables \((i,u)\), is
\[
      \Theta(i,u):\equiv I(i,\Run(i,u),u).
\]
It is an equation to one and hence quantifier free.  For the successor step,
work in the shifted context
\([\Theta(i,u),G(u)]\).  Instantiate \(\mathsf{step}\) with the terms
\[
     i,\\ \Run(i,u),\\ u.
\]
Eliminate the two implications using the guard and the induction hypothesis.
This yields
\[
 I(Si,\operatorname{step}(i,\Run(i,u),u),u).
\]
The primitive-recursion successor equation gives
\[
 \Run(Si,u)=\operatorname{step}(i,\Run(i,u),u).
\]
Transport in the state argument, reversing the equality when necessary with
\code{Prf.eqSymm}, to obtain
\begin{equation}
 [\Theta(i,u),G(u)]\vdash\Theta(Si,u).
 \label{eq:invstep}
\end{equation}

\paragraph{Step 3: invoke quantifier-free induction.}
Apply \code{PRA.qfInd} to \(\Theta\), using \eqref{eq:invbase} and
\eqref{eq:invstep}.  The result is
\begin{equation}
       [G(u)]\vdash \forall i\,I(i,\Run(i,u),u).
       \label{eq:allinv}
\end{equation}
The context shifting required by the induction rule is exactly the shifting
already built into \code{Prf.allI} and \code{PRA.qfInd}; no informal
renaming convention is needed.

\paragraph{Step 4: specialize at the computed bound.}
Use universal elimination in \eqref{eq:allinv} with the term
\(\operatorname{bound}(u)\).  We obtain
\begin{equation}
 [G(u)]\vdash
 I(\operatorname{bound}(u),
   \Run(\operatorname{bound}(u),u),u).
 \label{eq:boundinv}
\end{equation}

\paragraph{Step 5: apply the final contract.}
Instantiate \(\mathsf{final}\) with
\[
 s=\Run(\operatorname{bound}(u),u)
 \quad\text{and}\quad u=u.
\]
Weaken to the guard context and eliminate the two implications using the guard
and \eqref{eq:boundinv}.  This gives
\[
  [G(u)]\vdash
  A(\Run(\operatorname{bound}(u),u),u).
\]
Discharge \(G(u)\) by implication introduction and then close \(u\) by
universal introduction.  This is exactly \eqref{eq:traceresult}.

\subsection{Why the construction is primitive recursive on proofs}

The derivation is assembled from a fixed finite pattern of constructors plus
the three input derivations.  It performs no search through proofs.  Term and
formula substitutions are structural recursive functions already defined in
the checked library.  Hence, at the meta-level, \code{traceReflection} is an
ordinary total Lean definition.

\subsection{Multiple state components and certificate checks}

A single natural state can encode a tuple containing:
\begin{itemize}
\item program variables;
\item a current query-table index;
\item accumulated Boolean validity bits;
\item an error numerator and denominator;
\item a lexicographic descent measure;
\item a bounded certificate cursor.
\end{itemize}
The invariant remains a single Boolean equation checking all components.  A
large conjunction should be compiled as a Boolean fold, not represented by a
formula with a variable number of logical connectives.

\subsection{Nested loops}

Nested bounded loops may be flattened into one trace by pairing loop counters,
or handled compositionally.  Flattening is useful when one global invariant is
simpler than a hierarchy of intermediate theorems.  Composition is useful when
local loop proofs are reusable.  Both routes stay inside quantifier-free
induction because every loop is bounded and each local invariant has been
Booleanized.

\subsection{Termination by a ranked measure}

Some algorithms are described as ``repeat until done'' but come with a
primitive-recursive rank \(r(s)\) that strictly decreases when the terminal
predicate fails.  If an initial bound \(R(u)\) on the rank is supplied, compile
the algorithm to exactly \(R(u)\) steps, making terminal states stutter.  The
invariant records
\[
 \operatorname{done}(s_i)=1
 \quad\text{or}\quad r(s_{i+1})<r(s_i).
\]
A bounded induction proves that a nonterminal state cannot persist beyond
\(R(u)\) decreases.  Thus ranked descent is a special case of trace reflection,
not an additional induction principle.

\subsection{Lean realization}

The file \code{ObjectTraceCompiler.lean} contains arity-correct builders
\code{app1}, \code{app2}, and \code{app3}; the trace-system structure; all five
formula builders; the quantifier-free lemma; the run equations; and exact
signatures for four milestones:
\begin{enumerate}
\item \code{trace_base};
\item \code{trace_step};
\item \code{trace_all};
\item \code{trace_final}.
\end{enumerate}
The main definition \code{traceReflection} is then a short composition of
those milestones.  The remaining \code{sorry} obligations are substitution
normalizations and equality transports, not missing mathematical ideas.  The
last theorem
\[
 \code{traceReflectionNative}
\]
is already one application of the checked \code{pra_iff_native} bridge once
\code{traceReflection} is complete.

\subsection{Consequence for the proof packets}

Any packet whose main proof has been reduced to a fixed bounded algorithm with
local iPRA derivations can use \cref{thm:tracereflection}.  It no longer needs a
custom meta-argument that the algorithm ``ought to be formalizable.''  The
packet must provide exactly three kinds of evidence:
\begin{enumerate}
\item primitive-recursive symbols for the state transition and bound;
\item a quantifier-free invariant strong enough for the theorem;
\item closed object-language derivations of base, preservation, and final
soundness.
\end{enumerate}
This is the smallest reusable bridge between finite mathematics and the
checked logical backend.
